\clearpage
\section*{September 24, 2026: Non-optimizers and a heterogeneous jump}
\textit{Agent-drafted research entry.}

Jacob asked to try both standard ways of letting some parents stay above 100:
a share of non-optimizers, which he doubted for projects this large, and
parents that differ in how much the threshold costs them. Both are estimated
by the likelihood of the previous entry, with its unexplained share
$\varepsilon$, in \path{fit_heterogeneity.R}.

\begin{itemize}
\item \textbf{Non-optimizers.} A share $\pi$ of parents keeps its 2019--2022
outcome, as if the threshold did not reach it (Kleven and Waseem 2013). They
lose no units.
\item \textbf{Heterogeneous jump.} Each parent's jump is lognormal with median
$\kappa$ and log standard deviation $s$, so some parents barely feel the wage
rule and others always avoid 100. The distribution is placed on the $\kappa$
grid, extended to 5, where every parent of at most 300 units avoids 100; its
prediction is then a weighted average of predictions at single jumps.
\end{itemize}

Both nest the main model ($\pi=0$, $s=0$), and a third model combines them.
The single jump reproduces the likelihood estimate exactly.

\begin{center}
\small
\begin{tabular}{lrrrr}
\hline
 & Single jump & Non-optimizers & Heterogeneous jump & Both \\
\hline
Jump $\kappa$ (median) & 0.12 & 0.14 & 0.18 & 0.18 \\
Spread $s$ & -- & -- & 2 & 2 \\
Non-optimizer share $\pi$ & -- & 0.15 & -- & 0 \\
Kink $\tau$ & 0.05 & 0.05 & 0 & 0 \\
Cost growth $\gamma$ & 1.5 & 1.5 & 2 & 2 \\
Mean splitting cost $\sigma$ & 0.32 & 0.25 & 0.20 & 0.20 \\
Unexplained share $\varepsilon$ & 0.03 & 0.01 & 0.015 & 0.015 \\
Log-likelihood & $-592.3$ & $-586.3$ & $-584.5$ & $-584.5$ \\
AIC & 1,194.5 & 1,184.6 & 1,181.1 & 1,183.1 \\
Units lost & 1,130 & 930 & 761 & 761 \\
\hline
\end{tabular}
\end{center}

\paragraph{What the data prefer.} Both departures improve on the single jump:
the likelihood-ratio statistics are 11.9 for non-optimizers and 15.4 for the
heterogeneous jump, each one parameter on the boundary ($p=0.0003$ and
$p=0.00004$). The heterogeneous jump fits better with the same number of
parameters, and once it is allowed, the combined model sets $\pi$ to zero.
Against non-optimizers alone, adding the spread is significant ($p=0.03$).

The reason is the region just above 100. Non-optimizers keep their historical
size, so a 15 percent share puts 0.4 percent of recent parents at 100--104
units in one building, where none are observed. A parent with a small jump
still moves from 102 to 99, because that costs it almost nothing, so the
heterogeneous jump leaves that region empty while keeping parents of 105--149
units: 3.9 percent against 4.7 observed, from 2.7 under the single jump. It
also fits pairs of 199--300 units (0.5 against 0.7 observed, from 1.7) and
$99+99$ (3.4 against 3.2, from 2.8). It fits the spike at 99 worse (10.9
against 12.9, from 13.0). Jacob's doubt about non-optimizers is borne out in
this sense: the data favor parents who respond to different costs over
parents who do not respond.

\paragraph{The estimated spread.} With median 0.18 and $s=2$, the jump is
below 0.01 for about 6 percent of parents on the grid, below 0.05 for about a
quarter, and above 1 for about 17 percent, so the distribution is close to a
mix of nearly unaffected parents and parents who will not cross 100 at any
size up to 300. The likelihood falls on both sides of $s=2$ ($-584.7$ at 2.5,
$-586.0$ at 4). With heterogeneity the kink goes to zero: the kink had been
standing in for parents who stay above 100.

\paragraph{Units.} Units lost fall from 1,130 under the single jump to 761
with a heterogeneous jump, within the bootstrap interval of the single-jump
model (337--2,981). Parents who stay above 100 are now parents with small
jumps, who lose little, rather than parents trimmed by a common kink. These
are point estimates; the models with $\pi$ and $s$ have not been
bootstrapped.

\paragraph{Next.} The heterogeneous jump is the natural place for Jacob's
reserved wage test: if the jump reflects the gap between required and
prevailing construction wages, parents in areas that already paid more should
draw smaller jumps.

\subsection*{September 24 follow-up: a burden scaled by developer}

Jacob asked whether the heterogeneous jump forces the kink to zero. It
nearly does: with a common kink, the best fit has none, $\tau=0.02$ costs 0.7
log points and $\tau=0.05$ costs 2.0, just outside the 95 percent
likelihood-ratio interval. The two explain the same fact, that single
buildings of 150--300 units fell from about 21 to 8 percent of parents: the
kink by making large buildings costlier, the jump distribution by giving about
17 percent of parents a jump too large to cross at any size.

The wage rule, however, raises the cost of every construction hour, so a
parent that feels it little should feel both the jump and the kink little. The
scaled burden multiplies each parent's jump and kink by one lognormal scale
with median 1, so $\kappa$ and $\tau$ describe the median parent. Its best
point is again no kink, identical to the heterogeneous jump (log-likelihood
$-584.5$, 761 units lost), but the kink is now much less costly:

\begin{center}
\small
\begin{tabular}{rrrrr}
\hline
Kink-to-jump ratio & Median $\kappa$ & Median $\tau$ & Log-likelihood below best & Units lost \\
\hline
0 & 0.18 & 0 & 0 & 761 \\
0.2 & 0.12 & 0.024 & 0.44 & 826 \\
0.5 & 0.10 & 0.05 & 1.28 & 963 \\
0.75 & 0.10 & 0.075 & 1.94 & 1,142 \\
1 & 0.08 & 0.08 & 2.63 & 1,127 \\
\hline
\end{tabular}
\end{center}

A kink up to half the jump, including the single-jump estimate's
$\tau=0.05$, is inside the 95 percent interval, and units lost range from 761
to 963 within it. The data allow heterogeneity and a kink together; they
cannot tell whether the kink is zero or moderate. The scaled burden is the
natural home for the reserved wage test, which would make the scale depend on
the local wage gap.
